\cleardoublepage

\section{绪论}

\subsection{研究背景与意义}
积雪深度不仅是陆地水文和冰冻圈能量平衡的重要参数，也是海冰厚度反演中的关键输入量。对于北极海冰区域，积雪会影响海冰表面反照率、热传导过程以及卫星测高反演海冰厚度的精度，因此获取高空间分辨率、可持续观测的雪深数据，对冰冻圈变化监测和海冰厚度估计具有重要意义\citep{IPCC2019SROCC,Cohen2014}。

传统雪深观测主要包括地面站点、样线测量、浮标和机载雪雷达等。此类数据精度较高，但受限于观测成本和极地环境，难以提供连续的大范围覆盖。被动微波遥感具有长期序列和大范围覆盖优势，但空间分辨率较低，且容易受混合像元、雪湿度、冰型和地表状态影响。基于激光或雷达测高的雪深估计通常依赖雪期与无雪期高程差，或依赖外部数字高程模型、雷达测高等辅助数据，因此在数据匹配和表面参考基准构建方面仍存在限制\citep{DeschampsBerger2023_SnowDepth_ICESat2}。

ICESat-2 搭载的先进地形激光高度计 ATLAS 采用 532\,nm 光子计数激光雷达，能够提供光子级高程观测，为雪深探测提供新的技术条件\citep{ICESat2_GSFC,Smith2019_ICESat2_HeightAlgorithm}。与传统仅利用雪面几何高程信息的思路不同，光子计数激光雷达回波中由雪层内部多次散射形成的次表层长尾信号，记录了光子在雪介质中的传播路径统计特征，这为不依赖无雪期参考高程的直接雪深反演提供了新的信息来源\citep{Hu2022,Lu2022_DerivingSnowDepth}。因此，本文关注的核心问题是：能否利用 ICESat-2 光子计数激光雷达回波剖面中的次表层多次散射信号，在不依赖无雪期参考高程的条件下估计雪深，并通过 OIB 机载雪深数据对反演结果进行验证。

\subsection{相关研究与问题分析}
现有基于 ICESat-2 或激光测高数据的雪深探测方法主要可分为两类。第一类是高程差分方法，即通过比较积雪期与无雪期地表高程，或结合外部数字高程模型获得积雪厚度。这类方法物理含义直接，但依赖高质量参考高程，并且对时空配准精度要求较高。第二类是直接利用光子回波剖面的方法，即从雪层内部多次散射形成的次表层信号中提取路径长度分布信息，进而反演雪深。该类方法减少了对无雪期参考高程的依赖，是近年星载光子计数激光雷达雪深探测中的重要发展方向\citep{DeschampsBerger2023_SnowDepth_ICESat2,Lu2022_DerivingSnowDepth}。

Hu 等基于辐射传输理论和蒙特卡洛模拟，指出雪层内部多次散射形成的光子路径长度分布可近似服从 Gamma 分布，并建立了一阶、二阶和三阶路径长度统计矩与雪深之间的关系\citep{Hu2022}。Lu 等进一步将该方法应用于 ICESat-2 ATL03 光子数据，并分析了 afterpulse、表面粗糙度、大气前向散射、太阳背景噪声和探测器暗计数等误差来源\citep{Lu2022Unc,Lu2022_DerivingSnowDepth}。上述研究为利用 ICESat-2 多次散射信号直接反演雪深提供了理论和实验基础。

然而，将该方法应用于实际 ATL03 数据时仍面临若干问题。首先，ATL03 光子数据噪声较强，表面定位、剖面构建和背景扣除会直接影响次表层长尾信号的稳定性。其次，ICESat-2 回波剖面会受到仪器系统响应和 afterpulse 伪影影响，若不进行校正，容易造成雪深反演偏差。最后，反演结果需要与机载雪深观测进行时空配对验证，但 ICESat-2 与 OIB 在足迹尺度、轨迹位置和观测时间上并不完全一致，这会影响误差统计和结果解释。

因此，本文的重点不是重新建立完整的辐射传输理论，而是在已有多次散射雪深反演理论基础上，构建适用于 ATL03 光子数据的处理与反演流程，并通过 OIB 机载雪深数据评估该方法在北极海冰区域的反演表现。

\subsection{本文研究内容}
针对光子计数激光雷达雪深直接反演中表面定位、剖面构建、系统响应校正和结果验证等问题，本文以 ICESat-2/ATLAS ATL03 光子数据为基础，结合已有多次散射路径长度统计理论，开展雪深反演流程构建与实验验证。本文的研究重点并不是重新提出新的辐射传输理论，而是将已有理论方法落实到实际 ATL03 数据处理过程中，并分析其在北极海冰区域的反演效果。

本文主要开展以下三个方面的工作。

第一，构建 ATL03 光子数据预处理与后向散射剖面生成流程。针对 ICESat-2 强束光子数据，提取光子高程、时间、经纬度和背景噪声等信息，并通过空间筛选、高程粗筛、主表面识别和表面对齐等步骤，形成用于雪深反演的垂向后向散射剖面。在此基础上，通过滑动窗口聚合和背景噪声扣除，提高次表层多次散射信号的稳定性。

第二，实现基于多次散射路径长度分布的雪深反演方法。依据已有研究中路径长度统计矩与雪深之间的关系，对表面对齐后的回波剖面进行系统响应去卷积，削弱 afterpulse 和仪器脉冲拖尾对浅层信号的影响。随后，对校正后的剖面进行吸收校正和路径长度统计矩计算，并通过迭代方式更新吸收系数和散射相关参数，最终获得沿轨雪深反演结果。

第三，开展 ICESat-2 反演雪深与 OIB 机载雪深数据的对比验证。选取 2019 年 4 月北极海冰区域的典型重叠观测日期，将 ICESat-2 反演结果与 Operation IceBridge 机载雪深观测进行空间配对，计算 Bias、RMSE、MAE 和相关系数等指标，评估反演结果的绝对误差、日尺度一致性和局地变化响应能力，并讨论空间配对尺度、噪声抑制和深雪振幅压缩等因素对结果的影响。

通过上述工作，本文希望验证基于 ICESat-2 多次散射回波剖面的雪深反演流程在实际光子数据中的可行性，并为后续扩大样本区域、优化参数设置以及引入更多外部验证数据提供基础。

\subsection{论文结构安排}
全文共分为五章，参考文献和附录置于正文之后，不作为正文编号章节。

第一章为绪论，介绍雪深探测的研究背景与意义，梳理基于遥感和光子计数激光雷达的雪深反演研究进展，分析现有方法在数据处理、系统响应校正和结果验证方面存在的问题，并明确本文的研究内容。

第二章介绍基于光子计数激光雷达多次散射信息的雪深反演方法。该章首先说明 ICESat-2 ATL03 光子数据的基本特点和总体处理思路，然后介绍光子预处理、主表面识别、后向散射剖面构建、系统响应去卷积以及基于路径长度统计矩的雪深反演原理。

第三章介绍反演程序实现与关键参数设置。该章从计算流程角度说明 ATL03 光子读取、场景筛选、剖面构建、噪声抑制、去卷积和迭代反演的实现方式，并结合典型日期样本展示中间处理结果，分析关键参数对剖面质量和反演稳定性的影响。

第四章开展实验数据与结果分析。该章以 2019 年 4 月北极海冰区域 ICESat-2 与 Operation IceBridge 重叠观测数据为基础，进行空间配对和定量验证，计算 Bias、RMSE、MAE 和相关系数等指标，并从均值一致性、逐点相关性、振幅压缩和误差来源等方面讨论反演结果。

第五章为总结与展望，总结本文完成的主要工作和实验结论，指出当前方法在深雪区响应、局地变化保持和多源验证方面的不足，并对后续扩大数据范围、优化参数设置和引入更多验证数据提出展望。
